
\section{Summary of Results}

In this dissertation, a broad study of neutral meson production in proton-proton collisions at $\sqrt{s} = 200$~GeV was presented using \ac{sPHENIX} data. The analysis focused on the measurement of $\pi^{0}$ and $\eta$ mesons via two photon decay channels and reconstructed within the electromagnetic calorimeter. Furthermore, the implementation of absolute calibration of \ac{EMCal} was also discussed. The Ohio University group was one of the primary group responsible for the development and implementation of \ac{EMCal} absolute calibration within \ac{sPHENIX}. 


Invariant differential cross sections for $\pi^{0}$ and $\eta$ mesons were measured as functions of transverse momentum ($p_T$) over a broad kinematic range. The results were obtained using a combination of minimum bias and photon-triggered data, enabling access to both low and high $p_T$ regions. In addition, the $\eta/\pi^{0}$ ratio was evaluated as a function of transverse momentum, providing a complementary observable with reduced systematic uncertainties. Both the invariant cross section and $\eta/\pi^{0}$ ratio were evaluated for different acceptance regions, $|\eta| < 0.35$ and $|\eta| < 1.0$, both for comparison with \ac{PHENIX} and with full \ac{sPHENIX} pseudorapidity coverage.


Similarly, the analysis ($\eta$ to $\pi^0$ ratio) was extended to study the pseudorapidity dependence of meson production using Photon 4 \ac{GeV} triggered data. Measurements was performed in selected transverse momentum intervals, allowing for a detailed exploration of longitudinal aspects of hadron production. 
